\section{Background}

\subsection{PPO and Isaac Lab in Brief}

\textbf{Proximal Policy Optimisation (PPO)} is the policy-gradient algorithm used to train every policy in this report \cite{schulman2017}. PPO is an on-policy actor-critic method. The actor outputs an action distribution conditioned on the current observation; the critic estimates the value function used to compute the advantage of each action. At every update, PPO collects a fresh batch of trajectories from the current policy and maximises a clipped surrogate objective:

\[
L^{\mathrm{CLIP}}(\theta) = \hat{\mathbb{E}}_t \left[ \min\!\left( r_t(\theta)\, \hat{A}_t,\; \mathrm{clip}\!\left( r_t(\theta),\, 1 - \epsilon,\, 1 + \epsilon \right) \hat{A}_t \right) \right],
\]

where $r_t(\theta) = \pi_\theta(a_t \mid s_t) / \pi_{\theta_{\mathrm{old}}}(a_t \mid s_t)$ is the importance-sampling ratio between the new and old policies, $\hat{A}_t$ is the estimated advantage, and $\epsilon$ is the clip range. The clip prevents updates from moving the policy too far in a single iteration, which trades a small amount of optimisation speed for substantially improved stability. In practice this makes PPO robust to hyperparameter choices and reliable enough to train large simulated systems without manual tuning of the trust region.

\textbf{Isaac Lab} is the simulator used for every experiment in this report \cite{mittal2023}. It is a modular robot-learning framework layered over NVIDIA Isaac Sim that exposes thousands of parallel simulation environments on a single GPU. Each parallel environment runs an independent copy of the robot, the terrain, and the reward computation, and PPO collects gradient information from the union of all environments at every step. This parallelism is the reason a wide-range velocity-tracking policy can be trained on a single workstation in roughly an hour rather than over days \cite{rudin2022}. Isaac Lab also provides standard interfaces for observation composition, command sampling, termination conditions, and domain randomisation, all of which are used in this work.

\subsection{Velocity-Command Curricula}

Both curriculum rules studied in this report operate on the same input: a per-task tracking-reward signal computed over recently observed trajectories. They differ in how that signal is consumed to produce the next sampling distribution over commands. This section presents each rule in the notation of its source paper.

\subsubsection{Box Adaptive Curriculum (Task-Specific)}

Margolis et al.\ (2022) consider a velocity command sampled at the start of each episode from a probability distribution over a two-dimensional command space spanned by the forward linear-velocity command $\mathbf{v}_x^{\mathrm{cmd}}$ and the yaw-rate command $\boldsymbol{\omega}_z^{\mathrm{cmd}}$. The command space is represented as a discrete grid with resolution $[0.5\ \mathrm{m/s},\ 0.5\ \mathrm{rad/s}]$ centered at $[0\ \mathrm{m/s},\ 0\ \mathrm{rad/s}]$. The sampling distribution at episode $k$ is denoted $p_{\mathbf{v}_x,\boldsymbol{\omega}_z}^{k}(\cdot, \cdot)$.

The \textbf{Box Adaptive} variant maintains the joint distribution as a product of independent marginals,

\[
p_{\mathbf{v}_x,\boldsymbol{\omega}_z}^{k}(\cdot, \cdot) = p_{\mathbf{v}_x}^{k}(\cdot)\, p_{\boldsymbol{\omega}_z}^{k}(\cdot),
\]

with the per-component update rules

\[
p_{\mathbf{v}_x}^{k+1}(\cdot) \leftarrow f_v\!\left( p_{\mathbf{v}_x}^{k}(\cdot),\, r_{v_x^{\mathrm{cmd}}} \right), \tag{8a}
\]

\[
p_{\boldsymbol{\omega}_z}^{k+1}(\cdot) \leftarrow f_\omega\!\left( p_{\boldsymbol{\omega}_z}^{k}(\cdot),\, r_{\omega_z^{\mathrm{cmd}}} \right), \tag{8b}
\]

where $r_{v_x^{\mathrm{cmd}}}$ and $r_{\omega_z^{\mathrm{cmd}}}$ are the velocity-tracking rewards on the most recent episode whose command was $\mathbf{v}_x^{\mathrm{cmd}},\, \boldsymbol{\omega}_z^{\mathrm{cmd}}$. The specific Box Adaptive update is

\[
p_{\mathbf{v}_x}^{k+1}(\mathbf{v}_x^n) \leftarrow
\begin{cases}
p_{\mathbf{v}_x}^{k}(\mathbf{v}_x^n) & r_{v_x^{\mathrm{cmd}}} < \gamma, \\
1 & \text{otherwise},
\end{cases} \tag{10a}
\]

\[
p_{\boldsymbol{\omega}_z}^{k+1}(\boldsymbol{\omega}_z^n) \leftarrow
\begin{cases}
p_{\boldsymbol{\omega}_z}^{k}(\boldsymbol{\omega}_z^n) & r_{\omega_z^{\mathrm{cmd}}} < \gamma, \\
1 & \text{otherwise},
\end{cases} \tag{10b}
\]

where $\gamma \in (0, 1)$ is a success threshold, and the neighbouring commands $\mathbf{v}_x^n,\, \boldsymbol{\omega}_z^n$ are the adjacent discretised cells on each axis:

\[
\mathbf{v}_x^n \in \{\mathbf{v}_x^{\mathrm{cmd}} - 0.5,\ \mathbf{v}_x^{\mathrm{cmd}} + 0.5\}, \qquad
\boldsymbol{\omega}_z^n \in \{\boldsymbol{\omega}_z^{\mathrm{cmd}} - 0.5,\ \boldsymbol{\omega}_z^{\mathrm{cmd}} + 0.5\}.
\]

The structural property of this rule is that probability mass is only added; it is never removed. Once a cell is unlocked, it remains in the sampling distribution for the rest of training. The curriculum therefore expands outward from the initialisation in a monotone fashion.

In this report, the Box Adaptive rule is implemented over a one-dimensional command axis (forward linear velocity only, with lateral velocity and yaw rate fixed at zero per \S1.2) and is referred to as the \textbf{task-specific} curriculum.

\subsubsection{LP-ACRL (Teacher-Guided)}

Li, Li, and Hutter (2026) formulate curriculum learning as an explicit task-sampling problem. The task space is denoted $\mathcal{T}$, and the curriculum is a sequence of task-sampling distributions

\[
\mathcal{C} = (c_0, c_1, \ldots, c_j, \ldots), \qquad c_j \in \Delta(\mathcal{T}), \tag{1}
\]

where $c_j$ is the task-sampling distribution at curriculum stage $j$. A task instance is denoted $\zeta \in \mathcal{T}$.

The agent's performance on task instance $\zeta$ is measured by the expected episodic reward when training under distribution $c_j$:

\[
R_{c_j}(\zeta) = \mathbb{E}_{\tau \sim c_j}\!\left[ R_\tau \right], \tag{5}
\]

where $R_\tau$ is the cumulative reward of a trajectory $\tau$ of length $H$.

\textbf{Learning progress} on a task instance is defined as the change in expected episodic reward between two consecutive curriculum stages:

\[
LP_{c_j}(\zeta) = R_{c_j}(\zeta) - R_{c_{j-1}}(\zeta). \tag{6}
\]

A task instance with positive $LP$ is improving; a task instance with $LP$ near zero has plateaued, either because it has been mastered or because it is currently unlearnable.

The task-sampling distribution at the next curriculum stage is then a softmax over learning progress:

\[
c_{j+1}(\zeta) = \frac{\exp\!\left( LP_{c_j}(\zeta) / \beta \right)}{\sum_{\zeta' \in \mathcal{T}} \exp\!\left( LP_{c_j}(\zeta') / \beta \right)}, \tag{7}
\]

where $\beta$ is a temperature parameter that controls the sharpness of the distribution. Small $\beta$ produces a near winner-take-all distribution; large $\beta$ produces a near-uniform distribution.

The structural property of this rule is that probability mass is \textit{reallocated} at every curriculum stage. When a task instance plateaus, its softmax weight collapses, and the freed mass is redistributed across whichever instances are currently improving fastest.

In this report, LP-ACRL is implemented on the same one-dimensional command axis as the Box Adaptive variant (forward linear velocity, 8 discrete task instances of width $0.5$ m/s) and is referred to as the \textbf{teacher-guided} curriculum.

\subsubsection{Structural Comparison}

The two rules differ in three properties that have direct consequences for the analysis in later chapters.

\begin{itemize}
  \item \textbf{Reversibility.} Box Adaptive adds probability mass to cells and never removes it. LP-ACRL reweights the entire distribution at every stage based on current learning progress.
  \item \textbf{Trigger.} Box Adaptive uses a fixed-threshold trigger ($\gamma$): a cell unlocks the moment its tracking reward crosses $\gamma$. LP-ACRL has no fixed threshold; sampling weight is allocated proportionally to recent improvement.
  \item \textbf{Dilution.} Once $k$ cells are unlocked under Box Adaptive, each unlocked cell receives roughly $1/k$ of the sampling budget regardless of which cells are still improving. Under LP-ACRL, sampling weight follows learning progress, so a plateaued cell shares minimal budget regardless of how many other cells are being sampled.
\end{itemize}

These three differences appear repeatedly in the analysis of Chapter 4.

\subsection{Reward Design for Wide-Range Velocity Tracking}

\subsubsection{Structure of a Locomotion Reward}

A locomotion reward for velocity tracking is typically a weighted sum of three families of terms \cite{rudin2022}.

\begin{itemize}
  \item \textbf{Tracking terms} that produce a positive scalar when the actual body velocity matches the commanded velocity. The standard form is an exponential kernel on velocity error.
  \item \textbf{Smoothness and effort penalties} that produce a negative scalar when the policy commands jerky or high-torque actions. Typical members of this family penalise joint velocity, joint acceleration, joint torque, and the change in the policy's output action between consecutive control steps.
  \item \textbf{Bonus terms} that reward specific gait-pattern features, most commonly a feet-air-time bonus that pays out when a foot is in swing phase for at least a threshold duration.
\end{itemize}

The total reward at step $t$ is

\[
r_t = \sum_{i} w_i\, r_i(\zeta_t),
\]

where $r_i$ is the $i$-th term and $w_i$ is its weight. Tracking terms have $w_i > 0$, penalty terms have $w_i < 0$, and bonus terms have $w_i > 0$. The relative magnitudes of these weights determine which solution the policy converges to.

\subsubsection{The Upstream Isaac Lab Go2 16-Term Reward}

The starting point of this work is the upstream Isaac Lab Go2 reward function published by Unitree Robotics (2024) \cite{unitree2024}, composed of 16 weighted terms. The two tracking terms are exponential kernels on linear and angular velocity error:

\[
r_{\mathrm{lin}} = \exp\!\left( -\, \frac{\| v_{xy}^{\mathrm{cmd}} - v_{xy} \|^2}{\sigma_{\mathrm{lin}}^{2}} \right), \qquad w_{\mathrm{lin}} = 1.5, \quad \sigma_{\mathrm{lin}} = 0.5\ \mathrm{m/s},
\]

\[
r_{\mathrm{ang}} = \exp\!\left( -\, \frac{(\omega_z^{\mathrm{cmd}} - \omega_z)^2}{\sigma_{\mathrm{ang}}^{2}} \right), \qquad w_{\mathrm{ang}} = 0.75, \quad \sigma_{\mathrm{ang}} = 0.25\ \mathrm{rad/s}.
\]

Both kernels saturate at $1$ when the actual velocity matches the command, so the maximum reward contribution from the two tracking terms combined is $w_{\mathrm{lin}} + w_{\mathrm{ang}} = 2.25$.

The smoothness and effort group is

\[
r_{\mathrm{smooth}} = -\,|w_{\dot{q}}|\, \|\dot{q}\|^2 \;-\; |w_{\ddot{q}}|\, \|\ddot{q}\|^2 \;-\; |w_\tau|\, \|\tau\|^2 \;-\; |w_{\Delta a}|\, \|a_t - a_{t-1}\|^2 \;\le\; 0,
\]

with all four weights $w_{\dot{q}},\, w_{\ddot{q}},\, w_\tau,\, w_{\Delta a}$ negative in the configuration, so that the leading sign of each term in $r_{\mathrm{smooth}}$ is negative.

The remaining ten terms cover orientation, height stability, contact bonuses (including feet-air-time), and termination penalties. The upstream values for all 16 weights, and the specific modifications adopted in this work, are reported in Chapter 3.

\subsection{Synthesis}

The three streams above set up the experimental setting of Chapter 3 and the comparisons of Chapter 4. PPO and Isaac Lab (\S2.1) supply the optimisation algorithm and the parallel simulator under which a curriculum comparison is feasible. The two curriculum rules (\S2.2) provide qualitatively different hypotheses about how training compute should be allocated across a wide command range: Box Adaptive grows the sampling support monotonically with measured per-bin success, while LP-ACRL reallocates probability mass continuously based on observed learning progress. The reward function (\S2.3) defines what succeeding on a per-bin task actually means in scalar terms, and its 16 weighted components are the surface on which both tracking quality and locomotion style are eventually scored.
